%% chapters/03-metodologi-umum.tex
%% Bab III — Metodologi Umum (Fondasi Bersama Jalur A dan B)
%% Sumber: writings/dissertation-outline.md §Bab III
%% Fase C — dibuat dari 04-metodologi.tex + 03-dasar-teori.tex

\chapter{Metodologi Umum}
\label{bab:metodologi-umum}
% Alias label agar referensi dari 04-metodologi.tex ke bab:dasar-teori tetap resolve
% selama masa transisi Fase C; alias ini dihapus pada Fase D.
\label{bab:dasar-teori}

Bab ini menyajikan fondasi bersama yang digunakan oleh Bab~IV (diagnostik)
dan Bab~V (prognostik). Fondasi tersebut mencakup kerangka konseptual dua
jalur penelitian yang saling melengkapi, spesifikasi empat dataset publik,
prosedur praproses sinyal dan ekstraksi fitur, metrik evaluasi bersama untuk
kedua jalur, serta infrastruktur komputasi dan protokol reproduksi.
Metodologi spesifik tiap jalur dikembangkan lebih lanjut di awal Bab~IV
(jalur diagnostik) dan Bab~V (jalur prognostik) masing-masing.

% =================================================================
\section{Kerangka Konseptual Terintegrasi}
\label{sec:bab3_kerangka}
% =================================================================

Penelitian ini disusun dalam dua jalur paralel yang mencerminkan dua moda
operasi utama sistem perawatan prediktif berbasis vibrasi:
\textbf{Jalur~A (Diagnostik)} dan \textbf{Jalur~B (Prognostik)}.
Kedua jalur berbagi fondasi data, praproses, dan infrastruktur yang diuraikan
pada bab ini, kemudian masing-masing mengembangkan metodologi dan hasilnya
pada Bab~IV dan Bab~V, sebelum akhirnya berkonvergensi pada kerangka
\emph{Perawatan Prediktif} (\emph{Predictive Maintenance}, PdM)
\emph{multi-tier} tiga lapis yang disintesiskan pada Bab~VI.

Jalur~A berfokus pada \textbf{identifikasi jenis dan keparahan kerusakan
bantalan} menggunakan dataset CWRU dengan pendekatan progresif dari
model berbasis kernel (SVM/LR), algoritma berbasis pohon keputusan
(DT/RF/XGBoost), hingga jaringan dalam (WDCNN). Setiap pendekatan
dilengkapi modul \emph{Explainable AI} (SHAP) yang sesuai dengan
arsitekturnya. Keluaran utama Jalur~A adalah
\emph{Fault Signature Maps} (FSM), yaitu representasi \emph{fingerprint}
per kelas kerusakan pada resolusi sinyal mentah 2.048 titik.

Di sisi lain, Jalur~B berfokus pada \textbf{estimasi sisa umur pakai}
(\emph{Remaining Useful Life}, RUL) menggunakan tiga dataset prognostik
(PHM2012, XJTU-SY, IMS) dengan tiga arsitektur \emph{backbone}:
Mamba-xLSTM-Net, N-BEATS-xLSTM-RUL, dan SparseGate-TCN-RUL.
Interpretabilitas Jalur~B dicapai melalui \emph{Top-k Sparse Autoencoder}
(SAE) yang memetakan representasi laten model ke frekuensi karakteristik
bantalan (BPFx) secara statistik terverifikasi.

Gambar~\ref{fig:bab3_kerangka_terintegrasi} mengilustrasikan alur keseluruhan
kedua jalur beserta titik konvergensi di Bab~VI.

\begin{figure}[ht]
\centering
\resizebox{\textwidth}{!}{%
\begin{tikzpicture}[
    node distance = 0.5cm and 1.0cm,
    ds/.style  = {rectangle, rounded corners=3pt, draw, thick, fill=gray!15,
                  align=center, minimum height=0.85cm, minimum width=2.6cm,
                  font=\small},
    ja/.style  = {rectangle, rounded corners=3pt, draw, thick, fill=blue!10,
                  align=center, minimum height=0.85cm, minimum width=2.8cm,
                  font=\small},
    jb/.style  = {rectangle, rounded corners=3pt, draw, thick, fill=green!10,
                  align=center, minimum height=0.85cm, minimum width=2.8cm,
                  font=\small},
    conv/.style = {rectangle, rounded corners=3pt, draw, thick, fill=orange!15,
                   align=center, minimum height=1.0cm, minimum width=4.0cm,
                   font=\small},
    arr/.style  = {-{Stealth[length=5pt]}, thick},
  ]

  %% --- JALUR A (baris atas) ---
  \node[ds] (cwru) {CWRU\\{\scriptsize 10 kelas, 48~kHz}};
  \node[ja, right=1.1cm of cwru]  (bench)
       {Kernel / Tree / Deep\\{\scriptsize SVM · RF · WDCNN}};
  \node[ja, right=1.0cm of bench] (shap)
       {SHAP\\{\scriptsize Kernel/Tree/Deep}};
  \node[ja, right=1.0cm of shap]  (fsm)
       {\emph{Fault Signature}\\{\scriptsize \emph{Maps} (FSM)}};

  \node[font=\scriptsize\bfseries, text=blue!70!black, above=0.2cm of cwru]
       {Jalur A (Diagnostik, Bab~IV)};

  \draw[arr] (cwru)  -- (bench);
  \draw[arr] (bench) -- (shap);
  \draw[arr] (shap)  -- (fsm);

  %% --- JALUR B (baris bawah) ---
  \node[ds, below=2.6cm of cwru] (progds)
       {PHM2012 · XJTU-SY\\{\scriptsize IMS}};
  \node[jb, right=1.1cm of progds] (rul)
       {\emph{Backbone} RUL\\{\scriptsize Mamba-xLSTM · TCN}};
  \node[jb, right=1.0cm of rul] (sae)
       {Top-$k$ SAE\\{\scriptsize $d_{\mathrm{lat}} = 1024$}};
  \node[jb, right=1.0cm of sae] (bpfx)
       {Pemetaan BPFx\\{\scriptsize Pearson + \emph{bootstrap}}};

  \node[font=\scriptsize\bfseries, text=green!60!black, below=0.2cm of progds]
       {Jalur B (Prognostik, Bab~V)};

  \draw[arr] (progds) -- (rul);
  \draw[arr] (rul)    -- node[midway, above, font=\scriptsize]{laten $\mathbf{h}$} (sae);
  \draw[arr] (sae)    -- (bpfx);

  %% --- Konvergensi Bab VI ---
  \node[conv, right=2.8cm of shap, yshift=-1.3cm] (konv)
       {Bab~VI: Kerangka PdM \emph{Multi-Tier}\\
        {\scriptsize Edge~IoT $\to$ Edge~Server $\to$ Cloud}\\
        {\scriptsize Interpretabilitas Dua-Lapis (FSM + SAE-BPFx)}};

  \draw[arr] (fsm.east)  -| (konv.north);
  \draw[arr] (bpfx.east) -| (konv.south);

\end{tikzpicture}%
}
\caption{Kerangka konseptual terintegrasi: Jalur~A (biru) menggunakan dataset
         CWRU untuk \emph{benchmark} lintas golongan algoritma
         (kernel/pohon keputusan/jaringan dalam) dan menghasilkan FSM sebagai
         XAI level sinyal; Jalur~B (hijau)
         menggunakan tiga dataset prognostik untuk estimasi RUL dan pemetaan
         BPFx melalui Top-$k$ SAE. Kedua jalur berkonvergensi pada kerangka
         PdM \emph{multi-tier} di Bab~VI.}
\label{fig:bab3_kerangka_terintegrasi}
\end{figure}

Pemilihan dua jalur paralel (bukan satu jalur hirarkis tunggal) mencerminkan
realitas sistem PdM industri yang membutuhkan dua kemampuan yang berbeda
namun saling melengkapi: kemampuan diagnostik untuk pengambilan keputusan
jangka pendek (jenis dan lokasi kerusakan), dan kemampuan prognostik untuk
perencanaan perawatan jangka panjang (estimasi kapan komponen akan gagal).
Interpretabilitas pada kedua kemampuan diwujudkan melalui pendekatan yang
sesuai dengan skala waktu masing-masing: FSM untuk interpretabilitas level
sinyal dalam satu segmen pendek (Jalur~A), dan SAE-BPFx untuk
interpretabilitas konsep laten lintas ratusan rekaman degradasi (Jalur~B).

% =================================================================
\section{Dataset}
\label{sec:bab3_dataset}
% =================================================================

Empat dataset publik digunakan dalam penelitian ini: satu untuk Jalur~A
dan tiga untuk Jalur~B. Semua merupakan dataset \emph{benchmark} yang banyak
diacu dalam literatur perawatan prediktif bantalan.
Tabel~\ref{tab:bab3_dataset_ringkasan} merangkum karakteristik utama
keempat dataset.

\begin{table}[ht]
\centering
\caption{Ringkasan karakteristik empat dataset yang digunakan dalam penelitian.}
\label{tab:bab3_dataset_ringkasan}
\small
\begin{tabular}{@{}lllll@{}}
\toprule
\textbf{Dataset} & \textbf{Jalur} & \textbf{Bantalan/Kelas} & \textbf{Fs (kHz)} & \textbf{Tugas} \\
\midrule
CWRU    & A (Diagnostik)  & 1 \emph{rig}, 10 kelas & 48   & Klasifikasi kerusakan \\
PHM2012 & B (Prognostik)  & 17 bantalan            & 25,6 & Estimasi RUL \\
XJTU-SY & B (Prognostik)  & 10 bantalan            & 25,6 & Estimasi RUL \\
IMS     & B (Prognostik)  & 4 bantalan × 3 run     & 20   & Estimasi RUL \\
\bottomrule
\end{tabular}
\end{table}

\subsection{CWRU (Jalur A)}
\label{subsec:bab3_cwru}

Dataset CWRU (\emph{Case Western Reserve University Bearing Data Center})
merupakan dataset diagnostik bantalan yang paling banyak diacu dalam literatur
\citetitb{Smith2015CWRU}. \emph{Test rig} terdiri atas motor induksi 2~HP
yang menggerakkan poros dengan bantalan uji SKF~6205-2RS~JEM. Kerusakan
eksperimental diinduksi melalui \emph{electrical discharge machining} (EDM)
pada tiga komponen: \emph{inner race} (IR), \emph{outer race} (OR), dan
elemen gelinding (\emph{ball}). Akselerometer dipasang pada posisi
\emph{drive-end} (DE, 48~kHz) dan \emph{fan-end} (FE, 12~kHz); penelitian
ini menggunakan sinyal DE 48~kHz sebagai sinyal utama. Empat kondisi beban
diterapkan: 0, 1, 2, dan 3~HP. Dari kombinasi jenis, lokasi, dan tingkat
keparahan kerusakan, diperoleh 10 kelas kondisi sebagaimana dirinci pada
Tabel~\ref{tab:bab3_cwru_kelas}.

\begin{table}[ht]
\centering
\caption{Sepuluh kelas kondisi dataset CWRU, lokasi, tingkat keparahan,
         dan jumlah sampel setelah \emph{temporal split} 54/13/33\%.}
\label{tab:bab3_cwru_kelas}
\small
\begin{tabular}{@{}clccrr@{}}
\toprule
\textbf{Kelas} & \textbf{Kondisi} & \textbf{Lokasi} &
\textbf{Kerusakan (inci)} & \textbf{Latih} & \textbf{Uji} \\
\midrule
0 & Normal            & --  & --     & 124 & 75 \\
1 & \emph{Inner Race} & DE & 0,007 & 124 & 75 \\
2 & \emph{Inner Race} & DE & 0,014 & 124 & 75 \\
3 & \emph{Inner Race} & DE & 0,021 & 124 & 75 \\
4 & \emph{Outer Race} & DE & 0,007 & 124 & 75 \\
5 & \emph{Outer Race} & DE & 0,014 & 124 & 75 \\
6 & \emph{Outer Race} & DE & 0,021 & 124 & 75 \\
7 & \emph{Ball}       & DE & 0,007 & 124 & 75 \\
8 & \emph{Ball}       & DE & 0,014 & 124 & 75 \\
9 & \emph{Ball}       & DE & 0,021 & 124 & 75 \\
\midrule
  & \textbf{Total}    &    &       & \textbf{1.240} & \textbf{750} \\
\bottomrule
\end{tabular}
\end{table}

Segmentasi dilakukan dengan jendela non-tumpang-tindih sepanjang 2.048 titik.
Pembagian data mengikuti \textbf{pemisahan temporal} (bukan pemisahan acak)
untuk menghindari kebocoran informasi waktu: 54\% pertama dialokasikan untuk
\emph{training} (1.240 sampel), 13\% berikutnya untuk \emph{validation}
(310 sampel), dan 33\% terakhir untuk pengujian (750 sampel). Normalisasi
Z-score dihitung hanya dari \emph{training set} dan diterapkan ke dua set
lainnya tanpa modifikasi lebih lanjut.

Pada Bab~IV, dataset ini menjadi sumber utama untuk \emph{benchmark}
lintas golongan algoritma. Pada Bab~V, CWRU juga digunakan untuk validasi
\emph{cross-dataset} SAE; namun hanya tersedia $n = 10$ rekaman, sehingga
daya statistik terbatas ($p > 0{,}05$). Keterbatasan ini ditandai secara
eksplisit pada §V.8.

\subsection{PHM2012 / FEMTO-PRONOSTIA (Jalur B)}
\label{subsec:bab3_phm2012}

Dataset PHM2012 merupakan hasil kompetisi \emph{Prognostics and Health
Management} IEEE 2012 yang dikembangkan oleh FEMTO-ST Institute, Université
de Franche-Comté, menggunakan \emph{test rig} PRONOSTIA
\citetitb{Nectoux2012PRONOSTIA}. Karakteristik dataset dirangkum pada
Tabel~\ref{tab:bab3_phm2012}.

\begin{table}[ht]
\centering
\caption{Karakteristik dataset PHM2012 (FEMTO-PRONOSTIA).}
\label{tab:bab3_phm2012}
\small
\begin{tabular}{@{}ll@{}}
\toprule
\textbf{Atribut} & \textbf{Nilai} \\
\midrule
Bantalan uji        & NSK~6804 (\emph{deep groove ball bearing}) \\
Kondisi operasi     & 3 (1.800~rpm/4.000~N; 1.650~rpm/4.200~N; 1.500~rpm/5.000~N) \\
Total bantalan      & 17 \\
Sensor              & Akselerometer horizontal + vertikal (25,6~kHz) \\
Durasi rekaman      & 0,1 detik per rekaman; interval 10 detik \\
Sampel per rekaman  & 2.560 \\
Kriteria EOL        & RMS $\geq$ 20~g \\
Skema label RUL     & Linear 1,0 (sehat) $\to$ 0,0 (EOL) \\
\bottomrule
\end{tabular}
\end{table}

Skema pelabelan RUL yang digunakan dalam penelitian ini mengikuti pendekatan
\emph{piecewise}: nilai target dipertahankan pada 1,0 selama fase sehat
(sebelum titik onset degradasi $t_d$), kemudian turun secara linear ke 0,0
pada \emph{end-of-life} $t_f$:

\begin{equation}\label{eq:bab3_label_piecewise}
  r_t = \begin{cases}
    1{,}0,                            & t < t_d \\[4pt]
    \dfrac{t_f - t}{t_f - t_d},      & t_d \leq t \leq t_f
  \end{cases}
\end{equation}

\noindent Titik $t_d$ ditentukan secara otomatis melalui deteksi perubahan
signifikan pada nilai RMS menggunakan metode $3\sigma$ pada jendela geser.
Skema ini lebih realistis secara fisik dibandingkan skema linear penuh
karena mencerminkan fase sehat panjang yang diikuti fase degradasi akseleratif
yang relatif singkat.

Pembagian data mengikuti konvensi kompetisi PHM2012 secara berbasis bantalan
(\emph{bearing-wise}) sebagaimana dirinci pada
Tabel~\ref{tab:bab3_split_phm}. Pemisahan berbasis bantalan memastikan
bahwa tidak ada sampel dari bantalan yang sama tersebar ke dua set berbeda,
mencegah terjadinya kebocoran karakteristik degradasi individual.

\begin{table}[ht]
\centering
\caption{Pembagian set data PHM2012 (bearing-wise).}
\label{tab:bab3_split_phm}
\small
\begin{tabular}{@{}lll@{}}
\toprule
\textbf{Set} & \textbf{Bantalan} & \textbf{Jumlah} \\
\midrule
\emph{Training}   & 1\_1, 1\_2, 1\_4, 2\_1, 2\_3, 2\_5, 3\_1 & 7 \\
\emph{Validation} & 1\_5, 2\_2 & 2 \\
\emph{Test}       & 1\_3, 1\_6, 1\_7, 2\_4, 2\_6, 2\_7, 3\_2, 3\_3 & 8 \\
\bottomrule
\end{tabular}
\end{table}

Kompetisi PHM2012 mendefinisikan \emph{scoring function} asimetris yang
menjadi dasar PHM Score pada Subbab~\ref{subsec:bab3_metrik_prognostik}.
Pembagian yang mengikuti konvensi kompetisi memungkinkan perbandingan
langsung hasil penelitian ini dengan literatur \emph{benchmark} yang
dipublikasikan sebelumnya.

\subsection{XJTU-SY (Jalur B)}
\label{subsec:bab3_xjtusy}

Dataset XJTU-SY dikembangkan oleh Xi'an Jiaotong University bersama
Sumyoung Technology menggunakan \emph{test rig} akselerasi umur pakai
bantalan gelinding \citetitb{Wang2020XJTU}. Karakteristik dataset dirangkum
pada Tabel~\ref{tab:bab3_xjtusy}.

\begin{table}[ht]
\centering
\caption{Karakteristik dataset XJTU-SY.}
\label{tab:bab3_xjtusy}
\small
\begin{tabular}{@{}ll@{}}
\toprule
\textbf{Atribut} & \textbf{Nilai} \\
\midrule
Bantalan uji           & LDK UER204 (\emph{deep groove ball bearing}) \\
Kondisi operasi        & 3 (2.100~rpm/12~kN; 2.250~rpm/11~kN; 2.400~rpm/10~kN) \\
Total bantalan         & 15 (5 per kondisi) \\
Sensor                 & Akselerometer horizontal + vertikal (25,6~kHz) \\
Durasi rekaman         & 1,28 detik; interval 1 menit \\
Sampel per rekaman     & 32.768 \\
Mode kegagalan dominan & \emph{Outer-race spalling} (mayoritas bantalan) \\
Catatan                & Kondisi 3 tidak digunakan (keterbatasan akses data) \\
\bottomrule
\end{tabular}
\end{table}

Hanya kondisi 1 dan 2 yang digunakan, menghasilkan 10 bantalan aktif.
Pembagian set dirinci pada Tabel~\ref{tab:bab3_split_xjtu}.

\begin{table}[ht]
\centering
\caption{Pembagian set data XJTU-SY (bearing-wise, 2 kondisi aktif).}
\label{tab:bab3_split_xjtu}
\small
\begin{tabular}{@{}lll@{}}
\toprule
\textbf{Set} & \textbf{Bantalan} & \textbf{Jumlah} \\
\midrule
\emph{Training}   & 1\_1, 1\_2, 1\_3, 2\_1, 2\_2 & 5 \\
\emph{Validation} & 1\_4 & 1 \\
\emph{Test}       & 1\_5, 2\_3 & 2 \\
\bottomrule
\end{tabular}
\end{table}

Ketersediaan label mode kegagalan per bantalan pada XJTU-SY menjadi aset
penting untuk validasi interpretabilitas Jalur~B: apabila SAE memetakan
fitur ke BPFO, konsistensi dengan dominansi \emph{outer-race spalling} yang
terdokumentasi oleh pemilik dataset memberikan uji silang fisis yang kuat.

\subsection{IMS / \emph{Intelligent Maintenance Systems} (Jalur B)}
\label{subsec:bab3_ims}

Dataset IMS dikembangkan oleh \emph{Center for Intelligent Maintenance
Systems} (IMS), University of Cincinnati, menggunakan \emph{test rig} dengan
empat bantalan Rexnord ZA-2115 yang beroperasi secara bersamaan pada kecepatan
tetap 2.000~rpm di bawah beban radial 26,7~kN
\citetitb{Qiu2006WaveletBearing}. Karakteristik dataset dirangkum pada
Tabel~\ref{tab:bab3_ims}.

\begin{table}[ht]
\centering
\caption{Karakteristik dataset IMS.}
\label{tab:bab3_ims}
\small
\begin{tabular}{@{}ll@{}}
\toprule
\textbf{Atribut} & \textbf{Nilai} \\
\midrule
Bantalan uji       & Rexnord ZA-2115 (\emph{tapered roller bearing}) \\
Kecepatan poros    & 2.000~rpm (33,3~Hz) \\
Beban radial       & 26,7~kN \\
Sensor             & Akselerometer (20~kHz) \\
Sampel per rekaman & 20.480 (1 detik) \\
Interval akuisisi  & 10 menit \\
Run tersedia       & 3 (\emph{Run}~1: ${\sim}43$~jam; \emph{Run}~2: ${\sim}160$~jam;
                       \emph{Run}~3: ${\sim}46$~jam) \\
\bottomrule
\end{tabular}
\end{table}

Setiap \emph{run} memantau keempat bantalan (Bantalan~1–4) secara bersamaan.
Mode kegagalan yang terdokumentasi: pada \emph{Run}~1, Bantalan~3 mengalami
\emph{outer-race fault} dan Bantalan~4 mengalami kegagalan elemen gelinding
(\emph{rolling element}); pada \emph{Run}~2, Bantalan~1 mengalami
\emph{inner-race fault}. Pembagian data mengikuti protokol berbasis
\emph{run}: \emph{Run}~1 dan \emph{Run}~2 digunakan untuk \emph{training}
dan \emph{validation}, sedangkan \emph{Run}~3 untuk pengujian.

IMS dipilih sebagai dataset prognostik ketiga karena karakteristiknya yang
berbeda dari PHM2012 dan XJTU-SY: bantalan tipe \emph{tapered roller}
(bukan \emph{ball bearing}), operasi \emph{run-to-failure} multi-bantalan
simultan, dan distribusi mode kegagalan yang mencakup \emph{rolling element}
dan sangkar (\emph{cage}), tidak hanya \emph{outer race}. Keberagaman jenis
bantalan, mode kegagalan, dan karakteristik sinyal ini memperkuat uji
universalitas pada Bab~V.

% =================================================================
\section{Praproses Sinyal Bersama dan Ekstraksi Fitur HI 36-D}
\label{sec:bab3_praproses}
% =================================================================

Meskipun Jalur~A dan Jalur~B menggunakan representasi input yang berbeda
untuk model akhirnya, keduanya berbagi tiga tahap praproses awal:
(1) segmentasi jendela geser, (2) normalisasi Z-score berbasis
\emph{training set}, dan (3) ekstraksi fitur atau penggunaan langsung sinyal
tersegmentasi. Pemisahan pada tahap ketiga bergantung pada pendekatan
masing-masing model.

\subsection{Segmentasi Jendela Geser dan Normalisasi}
\label{subsec:bab3_segmentasi}

Segmentasi membagi rekaman sinyal panjang menjadi segmen pendek menggunakan
jendela geser dengan ukuran yang disesuaikan per dataset dan pendekatan model:
\begin{itemize}
  \item \textbf{CWRU (Jalur A):} jendela 2.048 titik, \emph{stride} penuh
        (tidak tumpang-tindih) untuk \emph{training} dan evaluasi. Panjang
        2.048 dipilih agar mencakup setidaknya satu periode impuls BPFx
        pada 48~kHz.
  \item \textbf{PHM2012, XJTU-SY, IMS (Jalur B):} jendela temporal berisi
        $L_{\rm win}$ rekaman berurutan (\emph{health indicator} per rekaman);
        konfigurasi $L_{\rm win}$ dan \emph{stride} diuraikan pada
        Bab~V~§V.1.
\end{itemize}

Normalisasi Z-score diterapkan per kanal menggunakan statistik dari
\emph{training set}:

\begin{equation}\label{eq:bab3_zscore}
  \tilde{x}_{\tau} = \frac{x_{\tau} - \mu_{\rm train}}{\sigma_{\rm train}},
\end{equation}

\noindent di mana $\mu_{\rm train}$ dan $\sigma_{\rm train}$ adalah mean dan
deviasi standar yang dihitung hanya dari \emph{training set}, kemudian
diterapkan secara seragam pada \emph{validation} dan \emph{test set} untuk
mencegah kebocoran informasi statistik.

\subsection{Health Indicator 36-D untuk Model Berbasis Kernel dan Pohon Keputusan}
\label{subsec:bab3_hi36d}

Model berbasis kernel (SVM/LR) dan algoritma berbasis pohon keputusan
(DT/RF/XGBoost) pada Bab~IV tidak beroperasi langsung pada sinyal mentah.
Dari setiap segmen sinyal CWRU, diekstraksi \textbf{vektor HI 36-D} yang
terdiri atas 18 fitur per kanal akselerometer (\emph{drive-end} dan
\emph{fan-end}), masing-masing mencakup:

\begin{itemize}
  \item \textbf{9 fitur domain waktu:} nilai RMS, kurtosis,
        \emph{crest factor}, \emph{peak-to-peak}, skewness, deviasi
        standar, mean, variansi, dan \emph{median absolute deviation} (MAD).
  \item \textbf{9 fitur domain frekuensi:} sentroid spektral, entropi
        spektral, dan 7 fitur energi pita yang mencakup pita frekuensi
        karakteristik bantalan (BPFO, BPFI, BSF, FTF) beserta harmonik
        terkait.
\end{itemize}

Dengan dua kanal (DE dan FE), total dimensi vektor fitur per segmen adalah
$18 \times 2 = 36$. Definisi matematis lengkap setiap fitur dirujuk ke
Lampiran~\ref{lamp:fitur}. Jendela FFT tipe \emph{Hann} digunakan
untuk mengurangi \emph{spectral leakage} pada fitur domain frekuensi.
Normalisasi Z-score tambahan diterapkan pada vektor HI 36-D setelah
ekstraksi, dengan statistik juga bersumber dari \emph{training set}.

\subsection{Representasi Input untuk Model Berbasis Jaringan Dalam}
\label{subsec:bab3_input_deep}

Dua jenis model jaringan dalam dalam penelitian ini menggunakan representasi
input yang berbeda dari HI 36-D:

\begin{enumerate}
  \item \textbf{WDCNN (Bab~IV, Jalur A):} menerima sinyal mentah 2.048 titik
        secara langsung ($\mathbf{x} \in \mathbb{R}^{2048}$), tanpa ekstraksi
        fitur manual. Kernel konvolusional berukuran 64 dengan \emph{stride}
        16 pada lapisan pertama berfungsi sebagai ekstraktor frekuensi adaptif
        otomatis \citetitb{Zhang2017WDCNN}.
  \item \textbf{Backbone RUL (Bab~V, Jalur B):} menerima sekuens vektor
        HI 5-D yang diekstraksi per rekaman melalui dekomposisi energi RMS
        lima pita frekuensi. Input berbentuk
        $\mathbf{X} \in \mathbb{R}^{L_{\rm win} \times 5}$ (sekuens $L_{\rm win}$
        rekaman, masing-masing 5 fitur energi). Prosedur praproses prognostik
        diuraikan pada Bab~V~§V.1.
\end{enumerate}

Perbedaan representasi ini mencerminkan perbedaan mendasar antara kedua tugas:
diagnostik mengklasifikasikan jenis kerusakan dari satu segmen waktu pendek
(2.048 titik $\approx$ 42,67 ms pada 48~kHz), sedangkan prognostik mengestimasi
degradasi dari sekuens ratusan rekaman yang mencakup seluruh lintasan degradasi.

% =================================================================
\section{Metrik Evaluasi Bersama}
\label{sec:bab3_metrik_bersama}
% =================================================================

Bagian ini mendefinisikan seluruh metrik evaluasi yang digunakan pada Bab~IV
dan Bab~V. Pengelompokan mengikuti tiga dimensi: kinerja klasifikasi diagnostik
(Jalur~A), kinerja prognostik (Jalur~B), dan validasi interpretabilitas
(FSM untuk Jalur~A; SAE-BPFx untuk Jalur~B).

\subsection{Metrik Diagnostik}
\label{subsec:bab3_metrik_diagnostik}

Kinerja klasifikasi dievaluasi dengan empat metrik. Akurasi mengukur proporsi
prediksi benar secara keseluruhan:

\begin{equation}\label{eq:bab3_akurasi}
  \text{Akurasi} = \frac{\sum_{k=1}^{K} \text{TP}_k}{N_{\rm test}},
\end{equation}

\noindent di mana $\text{TP}_k$ adalah jumlah prediksi benar pada kelas ke-$k$
dan $N_{\rm test}$ adalah total sampel uji. Macro-F1 menghitung rata-rata F1
per kelas tanpa pembobotan terhadap frekuensi kelas:

\begin{equation}\label{eq:bab3_macrof1}
  F1_{\rm macro} = \frac{1}{K} \sum_{k=1}^{K}
    \frac{2\, P_k R_k}{P_k + R_k},
\end{equation}

\noindent dengan $P_k$ (presisi) dan $R_k$ (rekall) per kelas ke-$k$.
Macro-F1 dipilih sebagai metrik primer karena tidak terbias terhadap kelas
mayoritas; penting untuk dataset yang seimbang per kelas seperti CWRU.
Weighted-F1 juga dilaporkan untuk kelengkapan. \emph{Confusion matrix}
$K \times K$ disajikan untuk analisis kesalahan per kelas.

Tiga metrik tambahan digunakan khusus untuk validasi \emph{Fault Signature
Maps} (FSM):

\begin{enumerate}
  \item \textbf{Indeks diskriminabilitas} ($D$): mengukur sejauh mana vektor
        FSM antar-kelas berbeda secara geometris, dinyatakan sebagai jarak
        \emph{pairwise} rata-rata relatif terhadap norma rata-rata:
        \begin{equation}\label{eq:bab3_discriminability}
          D = \frac{1}{\binom{C}{2}} \sum_{c_1 < c_2}
            \frac{\bigl\|\mathrm{FSM}_{c_1}^{\rm abs} - \mathrm{FSM}_{c_2}^{\rm abs}\bigr\|_2}
                 {\tfrac{1}{2}\bigl(\bigl\|\mathrm{FSM}_{c_1}^{\rm abs}\bigr\|_2
                  + \bigl\|\mathrm{FSM}_{c_2}^{\rm abs}\bigr\|_2\bigr)}.
        \end{equation}
        Nilai $D$ lebih besar berarti FSM lebih mudah dibedakan antar kelas.

  \item \textbf{Fraksi monotonik keparahan} ($M_c$): mengukur fraksi posisi
        sinyal di mana atribusi FSM meningkat secara monoton sejalan dengan
        tingkat keparahan kerusakan (0,007 $\to$ 0,014 $\to$ 0,021 inci):
        \begin{equation}\label{eq:bab3_monotonicity}
          M_c = \frac{|\{p : \mathrm{FSM}_c[p]_{0{,}014}
                          > \mathrm{FSM}_c[p]_{0{,}007}
                        \wedge \mathrm{FSM}_c[p]_{0{,}021}
                          > \mathrm{FSM}_c[p]_{0{,}014}\}|}{P},
        \end{equation}
        di mana $P = 2.048$ adalah total posisi sinyal dan
        $\mathrm{FSM}_c[p]_s$ adalah nilai FSM absolut kelas $c$ pada posisi
        $p$ untuk tingkat keparahan $s$.

  \item \textbf{Koefisien stabilitas \emph{split-half}} ($\rho$): mengukur
        konsistensi FSM antara dua subset acak separuh dari data uji,
        dihitung sebagai korelasi Pearson rata-rata lintas $C$ kelas:
        \begin{equation}\label{eq:bab3_stability}
          \rho = \frac{1}{C} \sum_{c=1}^{C}
            r_{\rm Pearson}\!\bigl(\mathrm{FSM}_c^{\rm abs}\big|_{\rm subset_1},\;
                                    \mathrm{FSM}_c^{\rm abs}\big|_{\rm subset_2}\bigr),
        \end{equation}
        di mana $\mathrm{subset}_1$ dan $\mathrm{subset}_2$ adalah dua bagian
        sama besar yang dipisah secara acak dari data uji. Nilai $\rho \geq 0{,}90$
        menunjukkan FSM yang sangat stabil antar subsampel.
\end{enumerate}

\subsection{Metrik Prognostik}
\label{subsec:bab3_metrik_prognostik}

Kinerja prediksi RUL dievaluasi dengan tiga metrik. \emph{Root Mean Squared
Error} (RMSE) digunakan sebagai metrik primer karena peka terhadap kesalahan
besar (kritis secara operasional):

\begin{equation}\label{eq:bab3_rmse}
  \text{RMSE} = \sqrt{\frac{1}{N} \sum_{i=1}^{N} (r_i - \hat{r}_i)^2},
\end{equation}

\noindent di mana $r_i$ adalah target RUL dan $\hat{r}_i$ adalah prediksi
model pada sampel ke-$i$. \emph{Mean Absolute Error} (MAE) sebagai metrik
pelengkap yang lebih \emph{robust} terhadap pencilan:

\begin{equation}\label{eq:bab3_mae}
  \text{MAE} = \frac{1}{N} \sum_{i=1}^{N} |r_i - \hat{r}_i|.
\end{equation}

PHM Score merupakan \emph{scoring function} asimetris dari kompetisi IEEE
PHM~2012 \citetitb{Nectoux2012PRONOSTIA} yang memberikan penalti lebih berat
pada prediksi terlambat (estimasi RUL terlalu besar) dibandingkan prediksi
lebih cepat:

\begin{align}
  E_i &= 100 \cdot \frac{r_i - \hat{r}_i}{r_i}, \label{eq:bab3_phm_ei} \\[4pt]
  A_i &= \begin{cases}
    \exp\!\Bigl(-\ln 0{,}5 \cdot \dfrac{E_i}{5}\Bigr),  & E_i \leq 0 \\[6pt]
    \exp\!\Bigl(+\ln 0{,}5 \cdot \dfrac{E_i}{20}\Bigr), & E_i > 0
  \end{cases}, \label{eq:bab3_phm_ai} \\[4pt]
  \text{PHM Score} &= \frac{1}{N} \sum_{i=1}^{N} A_i \in [0,\, 1].
  \label{eq:bab3_phm_score}
\end{align}

Nilai PHM Score lebih tinggi lebih baik; prediksi sempurna menghasilkan 1,0.
Karena PHM Score tidak memiliki konvensi resmi untuk XJTU-SY dan IMS, RMSE
digunakan sebagai satu-satunya metrik yang dapat dibandingkan secara adil
lintas ketiga dataset prognostik.

\subsection{Metrik Interpretabilitas Laten}
\label{subsec:bab3_metrik_interpretabilitas}

Tiga metrik digunakan untuk mengevaluasi korespondensi fitur SAE terhadap
frekuensi karakteristik bantalan (BPFx) pada Jalur~B.

\textbf{Hit-rate} $H_\varphi$ mengukur proporsi fitur SAE yang berkorelasi
moderat dengan amplitudo BPFx $\varphi$ ($\varphi \in
\{\text{BPFO, BPFI, BSF, FTF}\}$):

\begin{equation}\label{eq:bab3_hitrate}
  H_\varphi = \frac{\bigl|\{i : |r_{i,\varphi}| \geq 0{,}30\}\bigr|}
                   {d_{\rm lat}},
\end{equation}

\noindent di mana $r_{i,\varphi}$ adalah korelasi Pearson antara vektor
aktivasi fitur SAE ke-$i$ dan amplitudo \emph{envelope spectrum} pada pita
frekuensi $\varphi$, dan $d_{\rm lat} = 1024$ adalah dimensi laten SAE.
Ambang batas 0,30 mengikuti konvensi korelasi moderat yang lazim dalam
literatur analisis sinyal.

\textbf{Bootstrap 95\% CI} dihitung dengan $B = 1000$ resampel rekaman dari
dataset untuk memperoleh selang kepercayaan empiris pada $H_\varphi$:

\begin{equation}\label{eq:bab3_bootstrap_ci}
  \text{CI}_{95} = \bigl[H_\varphi^{(2{,}5\%)},\; H_\varphi^{(97{,}5\%)}\bigr],
\end{equation}

\noindent di mana persentil dihitung dari distribusi empiris $H_\varphi$
atas $B$ resampel dengan penggantian. Selang ini mencerminkan ketidakpastian
akibat variabilitas sampel lintas rekaman.

\textbf{\emph{Permutation p-value}} dihitung dengan $B = 1000$ permutasi
label rekaman terhadap nilai amplitudo BPFx; hit-rate pada setiap permutasi
dibandingkan terhadap hit-rate yang teramati (uji dua-sisi):

\begin{equation}\label{eq:bab3_permutation_pvalue}
  p = \frac{\bigl|\bigl\{b : H_\varphi^{(b)} \geq H_\varphi^{\rm obs}\bigr\}\bigr| + 1}
           {B + 1}.
\end{equation}

Nilai $p < 0{,}05$ menunjukkan bahwa korespondensi SAE–BPFx yang teramati
tidak dapat dijelaskan oleh kebetulan semata dan merupakan properti
\emph{emergent} dari representasi yang terlatih. Dua kelompok kontrol negatif juga digunakan:
model dengan inisialisasi acak Xavier (kontrol arsitektur) dan \emph{hidden
states} berupa Gaussian \emph{noise} (kontrol prosedur SAE). Rincian
eksperimen kontrol negatif dikembangkan pada Bab~V~§V.4.

% =================================================================
\section{Infrastruktur dan Protokol Reproduksi}
\label{sec:bab3_infrastruktur}
% =================================================================

Seluruh eksperimen dilakukan pada infrastruktur yang dirinci pada
Tabel~\ref{tab:bab3_infra}. Pelatihan \emph{single-card} (satu GPU per
\emph{run}) dipilih untuk memastikan reproduktibilitas hasil antar-\emph{seed}.

\begin{table}[ht]
\centering
\caption{Infrastruktur komputasi yang digunakan dalam seluruh eksperimen
         Jalur A dan Jalur B.}
\label{tab:bab3_infra}
\small
\begin{tabular}{@{}ll@{}}
\toprule
\textbf{Komponen} & \textbf{Spesifikasi} \\
\midrule
GPU                  & NVIDIA A40 (46~GB VRAM) \\
\emph{Framework} DL  & PyTorch Lightning 2.x \\
\emph{Framework} ML  & scikit-learn 1.x, XGBoost 2.x, SHAP \\
Presisi (Jalur B)    & \texttt{bf16-mixed} (AMP, Ampere+) \\
Presisi (Jalur A)    & FP32 \\
Python               & 3.11 (\texttt{virtualenv}) \\
\bottomrule
\end{tabular}
\end{table}

\textbf{Protokol multi-\emph{seed}.} Untuk menjamin rigor statistik pada
Jalur~B, setiap konfigurasi model × dataset dijalankan dengan tiga
\emph{seed} berbeda: 42, 43, dan 44. Total eksperimen pelatihan pada Bab~V:

\begin{equation}\label{eq:bab3_totalruns}
  \text{Total \emph{run}} = 3\ \text{model} \times 3\ \text{dataset}
                            \times 3\ \text{seed} = 27.
\end{equation}

Metrik kinerja dilaporkan sebagai rata-rata dan deviasi standar populasi
atas tiga \emph{seed}: $\overline{M} \pm \sigma_M$. Deviasi standar
mencerminkan sensitivitas hasil terhadap inisialisasi bobot, yang merupakan
sumber variabilitas utama pada \emph{backbone} berbasis memori sekuensial.
Pada Jalur~A (\emph{benchmark} kernel/pohon keputusan/jaringan dalam),
\emph{seed} ditetapkan dari masing-masing notebook reproduksi karena
pelatihan klasifikasi bersifat deterministik dan stabil pada kondisi yang sama.

\textbf{Strategi \emph{checkpoint}.} \emph{Checkpoint} terbaik disimpan
berdasarkan metrik validasi minimum (\texttt{val/rmse} untuk \emph{backbone}
RUL; $F1_{\rm macro}$ validasi untuk klasifikasi). \emph{Early stopping}
pada WDCNN menggunakan kesabaran 15 \emph{epoch}. \emph{Backbone} RUL
dijalankan selama 75 \emph{epoch} penuh tanpa \emph{early stopping} agar
kurva pelatihan antar model dapat dibandingkan secara adil dengan jumlah
langkah identik.

\textbf{Repositori dan reproduksi.} Seluruh kode pelatihan, konfigurasi
eksperimen, dan skrip evaluasi tersedia secara publik melalui repositori
yang tercantum pada Lampiran~\ref{lamp:kode}. Notebook
reproduksi tersedia untuk setiap karya sumber:
\texttt{Conference1\_Classification\_SVM\_LR.ipynb} (Bab~IV §IV.2),
\texttt{Conference2\_Classification\_Tree.ipynb} (Bab~IV §IV.3), dan
\texttt{Journal1\_Fault\_Signature\_Maps.ipynb} (Bab~IV §IV.4–IV.13).
Notebook untuk Bab~V akan tersedia setelah publikasi
\citetitb{TotoSuharto2025Journal2RUL}.
